\documentclass[11pt]{article}
\usepackage[margin=1in]{geometry}
\usepackage{amsmath,amssymb,amsthm,mathtools}
\usepackage[T1]{fontenc}
\usepackage{lmodern}
\usepackage{microtype}
\usepackage[colorlinks=true,allcolors=blue]{hyperref}
\newtheorem{theorem}{Theorem}
\newtheorem{lemma}[theorem]{Lemma}
\newtheorem{corollary}[theorem]{Corollary}
\newtheorem{proposition}[theorem]{Proposition}
\newcommand{\Rel}{\operatorname{Rel}}
\newcommand{\TV}{d_{\mathrm{TV}}}
\newcommand{\E}{\mathbb{E}}
\newcommand{\Prb}{\mathbb{P}}
\newcommand{\bits}{\operatorname{bits}}
\title{Approximating directed-cycle probability via two-terminal reliability}
\author{}
\date{Research draft --- November 16, 2026}

\begin{document}
\maketitle
\begin{abstract}
We reduce the probability that an independently edge-retained directed graph
contains a cycle to directed two-terminal reliability. One witness event per
edge and reliability-based conditional sampling give a classical union
estimator with controlled error under adaptive randomized queries. Combining
this reduction with the theorem of Feng, Fu, and Guo and a finite-bit
implementation of their simple algorithm gives an FPRAS for rational edge
probabilities. We bound oracle calls, random bits, and rational operand sizes.
A Lean development verifies the probabilistic reduction and composed resource
bounds, as well as induced-subgraph and parent-pointer correspondence laws.
Accuracy remains conditional on an external guarantee; its variance argument
and the machine-level runtime argument are not formalized.
\end{abstract}

\section{Introduction}
Let $G=(V,E)$ be a directed multigraph. Retain each edge $e$ independently
with rational probability $p_e\in[0,1]$, and write $Z$ for the probability
that the retained graph contains a directed cycle. An additive estimate
obtained by direct simulation need not give a relative approximation when
$Z$ is exponentially small. The existence of an FPRAS for $Z$ was asked by
Amarilli~\cite{question}, in connection with approximation on probabilistic
graphs~\cite{probgraphs}.

We reduce this problem to directed two-terminal reliability. The witnesses are edges with return paths: there are only $m=|E|$ such
events, even when the graph has exponentially many cycles. Their probabilities and conditional samples
can be obtained from reliability estimates. The union estimator itself is
classical~\cite{klm}; the contribution here is its application through this
oracle reduction, with explicit control of approximation, adaptive oracle
failures, and finite randomness.

Write $n=|V|$ and $m=|E|$. An FPRAS returns a relative $\varepsilon$-approximation
with probability at least $1-\delta$ in time polynomial in input length,
$\varepsilon^{-1}$, and $\log(\delta^{-1})$. Throughout, a reachability path
may have length zero. Thus a self-loop is
a cycle by itself. All oracle invocations use fresh independent randomness.
We use total variation in the convention
$\TV(P,Q)=\frac12\sum_x|P(x)-Q(x)|$.

\begin{theorem}[Oracle reduction]\label{thm:reduction}
Suppose directed two-terminal reliability on rational probabilities in
$(0,1]$ admits a randomized estimator for every rational accuracy $0<a<1$,
with relative error $a$, success probability at least $3/4$, rational output,
and worst-case bit complexity polynomial in input length and $a^{-1}$.
Then directed-cycle probability admits an FPRAS.
More precisely, for rational $0<\varepsilon\leq1/2$ and $0<\delta<1/2$, the
reduction uses accuracy
\[
 a=\frac{\varepsilon}{100m(m+1)},\qquad
 N=\left\lceil200m^2\varepsilon^{-2}
                \left\lceil\log_2(4/\delta)\right\rceil\right\rceil,
 \qquad T=m+2mN
\]
when the support is cyclic. It makes at most $T$ amplified queries, each
using $O(\log(T/\delta))$ base calls. Its output $\widehat Z$ satisfies
$\Prb[|\widehat Z-Z|>\varepsilon Z]\leq\delta$.
If $Z=0$, it returns zero deterministically.
\end{theorem}

Feng, Fu, and Guo give a reliability approximation algorithm for general
directed graphs~\cite{ffg}. Their arithmetic model does not directly imply
the bit-model hypothesis above. Section~\ref{sec:bits} supplies a
discretization and size analysis for their Algorithm S, yielding the
following consequence of their Theorem~9.

\begin{corollary}\label{cor:fpras}
Directed-cycle probability under independent rational edge retention has
an FPRAS with worst-case running time polynomial in the binary input length,
$\varepsilon^{-1}$, and $\log(\delta^{-1})$. Parallel edges and self-loops
are permitted.
\end{corollary}

\section{A polynomial family of witnesses}
Let $\mu$ be the product law on retained edge sets. For $e=(u,v)$ define
\[
 A_e=\{X:e\in X,\ v\text{ reaches }u\text{ in }X\setminus\{e\}\},
 \qquad w_e=\mu(A_e)=p_e\Rel_{G-e}(v,u).
\]
An outcome is cyclic if and only if it belongs to some $A_e$. Indeed, an
edge on a cycle has a return path after its deletion; conversely a return
path together with $e$ contains a cycle. Define
$t(X)=|\{e:X\in A_e\}|$ and $W=\sum_e w_e$. If $Z>0$, then
\begin{equation}\label{eq:witness}
 Z\leq W\leq mZ.
\end{equation}
Choose $e$ with probability $w_e/W$, then draw $X$ from $\mu(\cdot\mid A_e)$.
The resulting law $Q$ satisfies
\[
 Q(X)=\frac{t(X)\mu(X)}{W}.
\]
Set $Y(X)=1/t(X)$ when $t(X)>0$, and $Y(X)=0$ otherwise. Then $0\leq Y\leq1$
on every outcome, and
\begin{equation}\label{eq:mean}
 \theta=\E_QY=Z/W\geq1/m.
\end{equation}
Both membership in $A_e$ and $t(X)$ are computable by reachability tests.
The positive-probability support contains a cycle precisely when $Z>0$,
so the zero case is decidable by a deterministic graph search.

\section{Conditional sampling with approximate reliability}
To sample conditionally on $A_e$, fix $e$ present and condition the other
edges on the return-path event in $G-e$. Expose those edges in a fixed order.
At a feasible prefix, let $R_b$ denote residual reliability when the next
edge $f$ is fixed to $b\in\{0,1\}$. Its two conditional weights are
\[
 c_0=(1-p_f)R_0,\qquad c_1=p_fR_1.
\]
A zero weight can be detected exactly: test its probability factor and
reachability in the positive-probability support. No oracle call is needed
for a zero weight. A positive weight uses one reliability estimate. A
residual instance only pins original probabilities to zero or one; delete
zero-probability edges before invoking the oracle. Cases with equal
terminals have reliability one and may be answered directly.

\begin{lemma}\label{lem:sampling}
Suppose every oracle answer on a positive instance is within relative error
$a\leq1/2$. Replacing exact weights by these estimates changes a conditional
sample law by total variation at most $4ma$. Including approximate witness
selection changes the law $Q$ by at most $4(m+1)a$.
The same bounds hold for an always-accurate randomized oracle using fresh
randomness for each call.
\end{lemma}
\begin{proof}
For nonnegative weights $c_i$, replace each positive weight by
$\widehat c_i\in[(1-a)c_i,(1+a)c_i]$ and leave zero weights zero. For each
positive coordinate, the ratio of normalized probabilities lies between
$(1-a)/(1+a)$ and $(1+a)/(1-a)$. The $\ell_1$ distance is consequently at
most $2a/(1-a)\leq4a$, which also bounds total variation. At each shared
prefix couple the exact and approximate next choices with disagreement
probability at most $4a$. There are at most $m$ exposures, and their
disagreement probabilities add. Apply the same normalization bound once
more to witness selection. For randomized answers the bound holds for
every realization of the answers at a shared prefix, and hence also after
averaging. Fresh oracle randomness defines a fixed transition law at each
prefix.
\end{proof}

\section{The algorithm and its error analysis}\label{sec:accuracy}
After the support test, estimate all positive witness weights using the
parameters in Theorem~\ref{thm:reduction}. Put
$\widehat W=\sum_e\widehat w_e$. Generate $N$ independent sampling runs
conditional on these initial estimates, each choosing a witness according
to $\widehat w_e/\widehat W$ and exposing its remaining edges as above.
Return
\begin{equation}\label{eq:output}
 \widehat Z=\widehat W\frac1N\sum_{j=1}^NY(X_j).
\end{equation}
For completeness on failed-oracle executions, replace negative answers
by zero; if all weights at a choice vanish, choose a fixed alternative.
Use the definition $Y=0$ for an outcome with no witnesses. These conventions
define a terminating algorithm on every execution and do not change
any execution with accurate answers.

\begin{proof}[Proof of Theorem~\ref{thm:reduction}, accuracy]
First consider exact discrete choices and an oracle that always returns
$a$-accurate answers. Conditional on the initial estimates, the samples
are independent with common expected score $\widehat\theta$, and
Lemma~\ref{lem:sampling} gives
\[
 |\widehat\theta-\theta|\leq4(m+1)a=\frac{\varepsilon}{25m}.
\]
Hoeffding's inequality and the specified $N$ give
\[
 \Prb\left[\left|\frac1N\sum_jY(X_j)-\widehat\theta\right|
       >\frac{\varepsilon}{10m}\right]
 \leq2\exp\left(-\frac{N\varepsilon^2}{50m^2}\right)\leq\delta/4.
\]
The bound is uniform in the initial estimates, so it holds unconditionally.
By~\eqref{eq:mean}, the relative error of the score average is at most
$(1/25+1/10)\varepsilon<\varepsilon/4$ on this event's complement.
Also $|\widehat W/W-1|\leq a\leq\varepsilon/200$. Multiplication in
\eqref{eq:output} gives relative error at most
$\varepsilon/4+\varepsilon/200+\varepsilon^2/800<\varepsilon$.

Now use a randomized base oracle. Take the median of $k$ independent calls,
where $k=8\lceil\log_2(4T/\delta)\rceil+1$ is odd and at least
$8\ln(4T/\delta)$. The integer ceiling is computable by rational comparisons
and repeated doubling.
The chance that at least half the calls fail is at most $e^{-k/8}$ by
Hoeffding's inequality, and thus at most $\delta/(4T)$. On each fixed query
define a fictitious oracle that replaces an inaccurate median by the exact
reliability. This oracle is always accurate and uses fresh randomness.
Couple its execution to the real execution until the first replacement.
Conditional on any shared history, the next query is fixed and its failure
probability is at most $\delta/(4T)$. There are at most $T$ queries, so a
union bound gives disagreement probability at most $\delta/4$.
This argument never conditions the samples on all oracle calls succeeding.

Finally replace exact discrete choices by fixed-bit choices. For $h$
alternatives, draw $U$ uniformly from $\{0,\ldots,2^b-1\}$ and locate
$U/2^b$ in their half-open cumulative-probability intervals. Each interval
has mass error at most $2^{-b}$, so total variation is at most $h/2^b$.
One sampling run has at most one $m$-way choice and $m$ binary choices.
Choose
\[
 b=\left\lceil\log_2(12Nm/\delta)\right\rceil.
\]
Sequential coupling over all $N$ runs changes any event probability by
at most $3Nm/2^b\leq\delta/4$, including on failed-oracle histories.
The three failure budgets sum to at most $3\delta/4<\delta$.
The support test handles $m=0$ and all other zero instances before any
parameter involving division by $m$ is used.
\end{proof}

\section{Bit complexity and the concrete oracle}\label{sec:bits}
\subsection{Cost of the reduction}
Let $L$ bound the bit length of each numerator and denominator in the input,
including $\varepsilon$ and $\delta$. A base query has $n$ vertices, at most
$m$ edges, and parameter bit length $L'=O(L+\log(m+1))$. Let $R$ bound its
worst-case bit cost and $H$ bound the bit length of each output numerator
and denominator, taking maxima over the possible residual instances.
Both bounds are evaluated at accuracy $a$. A median is one of these outputs,
so amplification does not increase $H$.

A branch weight is a product of an input probability and an oracle answer;
its numerator and denominator have $O(L+H)$ bits. A sum of at most $m$ witness
weights can be represented using their product denominator, with
$O(m(L+H)+\log(m+1))$ bits. Choice comparisons use these unnormalized
weights: multiply by $U$ and $2^b$ and compare integer partial sums, without
forming real probabilities. For the score sum use the fixed denominator
$m!$, since every positive $t(X)$ lies in $\{1,\ldots,m\}$. The accumulator
then has $O(m\log(m+1)+\log(N+1))$ bits. Multiplying it by $\widehat W$ gives
the rational output without any uncontrolled growth of denominators.

Consequently all integers in this implementation have at most a constant
multiple of
\[
 K=m(L+H+\log(m+1))+\log(N+1)+b
\]
bits, apart from parameter arithmetic, which has polynomial cost in the
input length. With schoolbook arithmetic, a conservative bit-cost bound is
\begin{equation}\label{eq:cost}
 O\!\left(T\left[(n+m)\log(n+m+2)
       +k\bigl(R+\log(k+1)K^2\bigr)\right]\right),
\end{equation}
plus that parameter computation. Here graph searches, witness tests, and
choice scans total $O(T)$ searches or $O(T)$ arithmetic operations, and
sorting the $k$ oracle answers accounts for the median term. Edges and
vertices can be indexed with $O(\log(n+m+2))$ bits.
This proves the cost assertion in Theorem~\ref{thm:reduction}.
In particular no bound depends polynomially on $1/Z$.

\subsection{Finite-bit reliability}
For the mathematical corollary we implement the preprocessing of
Feng--Fu--Guo explicitly: return one for equal terminals, delete vertices
unreachable from the source, and return zero if the target disappears.
In the remaining graph $n\geq2$, $m\geq n-1$, and all probabilities are
positive. Their Algorithm S with
\[
 M=\left\lceil16n^3m/a^2\right\rceil
\]
has relative error at most $a$ except with probability $1/8$, by their
Theorem~9~\cite{ffg}. Algorithm S grows a resolved set $R$ from $\{s\}$.
Each resolved vertex stores an estimate $\widehat q_u$ and a pool of DFS
discovered sets; $M$ independent sampling iterations generate the estimate
and pool of one new vertex per round. We replace its choices by fair bits.

\begin{lemma}[Rational discretization]\label{lem:discrete}
For rational edge probabilities, Algorithm S at this value of $M$ admits
a finite-bit implementation with failure probability at most $1/4$,
polynomial worst-case bit complexity, and polynomial-length rational output.
\end{lemma}
\begin{proof}
In one sampling iteration the choices are a resolved vertex (at most $n$
alternatives), a pool entry (at most $M$), the first open outgoing cut edge
(at most $m$), and at most $m$ binary edge states. There are at most $n-1$
rounds and $M$ iterations per round. Thus the sum of numbers of alternatives
along any execution is at most
\[
 A=(n-1)M(n+M+3m).
\]
Use $b_S=\lceil\log_2(8A)\rceil$ bits per choice. The sequential coupling
argument from Section~\ref{sec:accuracy} bounds the change in any output
event by $A/2^{b_S}\leq1/8$.

Write $p_e=c_e/d_e$ and $B=\prod_ed_e$. Every probability that a particular
cut edge is the first open one, and every probability $\gamma_u$ that the
cut at $u$ has some open edge, is a multiple of $1/B$ in $[0,1]$.
After $r$ completed rounds, all current estimates can be written over the
common denominator $(MB)^r$, with magnitude at most $n^r$. This holds
initially; a new estimate is $\widehat\Phi J/M$, where
$\widehat\Phi=\sum_{u\in R}\widehat q_u\gamma_u$ and $0\leq J\leq M$.
The denominator gains only a factor $MB$ and the magnitude at most a factor
$n$. Existing estimates can be rescaled to the new denominator.

This invariant controls intermediate selection weights as well as stored
estimates. At the start of a round, products $\widehat q_u\gamma_u$ have
common denominator $(MB)^rB$, and their partial sums have only an additional
$O(\log n)$ numerator bits. The first-open-edge weights share denominator
$B$. Sampling a pool of size $\ell$ uses index
$\lfloor U\ell/2^{b_S}\rfloor$, avoiding an explicit scan through its
$\ell$ equal-weight intervals. Comparisons of unnormalized partial sums
implement all remaining choices using integers of
\[
 K_S=O\bigl(n(mL+\log(M+1)+\log(n+1))+b_S\bigr)
\]
bits. Snapshots are $n$-bit vertex sets, and there are at most $nM$ of them.
A direct implementation using scans, explicit snapshots, and common
denominators needs at most $O(nM(n+m)^2)$ elementary arithmetic or set
operations. Its bit cost is bounded conservatively by
$O(nM(n+m)^2K_S^2)$, in addition to preprocessing and parameter computation.
All these bounds are polynomial in input length and $a^{-1}$, on every
execution. Adding the discretization error to $1/8$ proves the lemma.
\end{proof}

Theorem~\ref{thm:reduction} and Lemma~\ref{lem:discrete} prove
Corollary~\ref{cor:fpras}. The fast reliability algorithm's advertised
unit-cost running time is not a running-time bound for this construction.

\section{Correspondence and formal verification}\label{sec:formal}
The companion Lean development transcribes Algorithm S without deleting
unreachable vertices. The following written lemma connects that choice to
the preprocessed algorithm. It also records the invariants needed to
interpret the transcription's totalized operations.

\begin{lemma}[Restriction to reachable vertices]\label{lem:restriction}
Fix $M\geq1$, a finite directed graph with positive edge probabilities,
distinct terminals $s,t$, and a fixed arc order. Suppose $t$ is reachable
from $s$, and let $U$ be the vertices reachable from $s$ in the underlying
graph. The exact transcription and Algorithm S on the graph induced by
$U$, using the restricted arc order and the same $M$, have the same output
law. Stopping a sampling DFS when its selected target is popped, replacing
its parent-path test by a Boolean flag, and imposing fuel $|V|+1$ on DFS
and $|V|$ on rounds do not change that law.
\end{lemma}
\begin{proof}
There is no edge from $U$ to $V\setminus U$. Initially the resolved set and
the sole snapshot are $\{s\}$. Inductively all resolved vertices and stored
snapshots lie in $U$: cut edges with resolved tails end in $U$, and DFS
cannot leave $U$. Restriction therefore preserves the ordered cut, its
maximizing edge, all $\gamma_u$, and all nonzero vertex-selection weights.
Vertices outside $U$ have zero selection weight. Couple the choices on
their common nonzero alternatives.

The identity $\widehat q_u>0$ if and only if the pool at $u$ is nonempty
holds for resolved vertices. It holds at initialization. A continuing round
has a positive maximizing edge product, hence $\widehat\Phi>0$, and its
new estimate $\widehat\Phi J/M$ is positive exactly when $J>0$ records are
stored. Thus a selected vertex always has a nonempty pool. The extra null
alternatives and out-of-range pool indices in the transcription have
probability zero. Uniform pool indices and uniform pool entries agree,
including multiplicities.

Choosing the first open cut edge $e_j$ with weight
$p_{e_j}\prod_{i<j}(1-p_{e_i})$, fixing earlier cut edges closed, and
revealing later ones independently gives exactly the cut's product law
conditioned on at least one open edge. Revealing the remaining states only
when examined preserves the joint law of everything the DFS observes.
The initial snapshot contains its starting vertex, so that vertex is not
pushed again. Every other vertex is marked discovered when first pushed
and is pushed at most once. Hence a DFS has at most $|U|$ pops.

Give the starting vertex flag true. When discovering $y$ from a vertex with
flag $f$, give it flag $f\wedge(y\notin R)$. By induction this flag is true
exactly when its tree path avoids resolved vertices other than the root.
The snapshot at the target is taken when it is popped, before scanning its
outgoing edges. Continuing beyond that pop cannot alter the recorded
snapshot or its eligibility, and cannot record the target a second time.
The early stop and flag therefore preserve each iteration's record law.
Prepending instead of appending records only permutes a pool, preserving
future uniform draws. Finally each continuing round resolves a new vertex,
so at most $|U|-1$ rounds occur. Both fuel bounds are sufficient. Induction
over rounds proves equality of the output laws.
\end{proof}

If $M$ is computed using the unreduced graph, it is at least the budget
computed on the reduced graph. This does not invoke a general monotonicity
principle for sample counts: the variance proof in Section~2.4 of
\cite{ffg} gives $\operatorname{Var}(\widehat q_t)\leq2n'^3m'q_t^2/M$
whenever $M\geq2n'^3m'$, where $n',m'$ are the reduced sizes.
The unreduced budget therefore gives failure at most $1/8$ at error $a$.
Equal terminals return one directly, including on a one-vertex or edgeless
graph. The larger-budget variance argument remains written mathematics;
the equal-terminal return is also proved in Lean.

\subsection{Verified guarantees and remaining boundary}
Lean proves several components of Lemma~\ref{lem:restriction}: invariance
of the entire output law under changes outside a closed set containing the
source; sufficient DFS and round fuel; pool positivity, size, and membership
invariants; uniform pool sampling with multiplicities and under permutations;
and the first-open conditional product law. It also proves that stopping
at the target pop preserves the law of a snapshot recorded before finishing
DFS, and that the Boolean path fold expresses avoidance of resolved vertices.

Two whole-program correspondence laws close the subtype and parent-pointer
steps. The induced graph has actual vertex and edge subtypes and inherits
the original arc list and first-open ranks, with the same $M$ and ambient
fuel bounds. Its output law equals that of the transcription
(\texttt{algS\_inducedClosed}). A separately defined program stores and updates
parent pointers, tests the path when the target is popped, and finishes
exploring after freezing the snapshot. For $M>0$ and probabilities in $[0,1]$,
its full output law equals that of the flag-based transcription
(\texttt{parentAlgS\_law}). The parent traversal also satisfies the fuel bound.
These theorems concern mathematical program semantics; identifying that
model with the published pseudocode still requires review. The external
variance proof, including its larger-budget application, is not formalized.
Accordingly, Lean assumes \texttt{AlgSAccurate} for every accuracy in $(0,1)$:
the exact transcribed Algorithm S has relative error at most that accuracy
with probability at least $7/8$. This is a hypothesis about real probabilities;
the binary-input corollary only requires rational probabilities and uses
explicit preprocessing.

The formal reduction uses $a=\varepsilon/(128m(m+1))$ and
$N=\lceil2048m^2/\varepsilon^2\rceil$, with Chebyshev's inequality followed
by median amplification of whole runs. Thus its constants and confidence
argument differ from Section~\ref{sec:accuracy}; its failed-oracle fallback
conventions also differ. Lean verifies the witness identity, sampling error,
adaptive-oracle coupling, rounding, and resource bounds for that program.

The final theorem \texttt{cycleFPRAS\_spec} concerns one call-free program
with every reliability query expanded and both median amplifications
included. It combines the accuracy claim with a bound on all fair bits
and the height of every final output. The edgeless case returns zero.
For $m\geq1$, write $N$ for the formal sample budget, $T=m+2mN$, $k$ for the inner amplification count,
$r$ for the outer count, and $b$ for the reduction's draw width. With
$M_j=\lceil16n^3j/a^2\rceil$, define
\[
 C_j=nM_j\bigl((n+1)+(M_j+1)+(j+1)+(n+1)2j\bigr),\qquad
 b_j=\lceil\log_2(8C_j)\rceil,
\]
where the formal natural ceiling logarithm returns zero at zero. The
proved total fair-bit bound is
\[
 brN(m+1)+rTk\sum_{j=0}^{m}b_jC_j.
\]
For input denominators at most $H\geq1$, put
\[
 H_O=1+\sum_{j=0}^{m}\max\{1,[(nj+1)M_jH^j]^n\}.
\]
Every output has numerator absolute value and denominator at most
$\max\{1,(2HH_O)^mN m!\}$. The sums cover all residual positive-edge
subgraphs. Each nonempty median is the lower sample median, hence an
actual input value; neither amplification increases height. Empty medians
return zero, and the total bounds include zero sample budgets and
edgeless or equal-terminal residual queries. These resource statements do
not assume successful oracle answers or the external accuracy hypothesis.
The operation charges are annotations; converting the resource bounds into
a machine-level runtime remains the written argument of Section~\ref{sec:bits}.

The development uses Lean 4.35.0 and a pinned mathlib dependency. Its
audit of 100 declarations permits only \texttt{propext},
\texttt{Classical.choice}, and \texttt{Quot.sound}; there are no admitted
proofs in the audited development. From the repository root, reproduce the
build and audit with:
\begin{quote}\ttfamily
python3 scripts/check\_formalization.py
\end{quote}
Exact enumeration comparing two Python ports is additional finite evidence;
it neither proves correspondence nor links Python execution to Lean.

\paragraph{Provenance.}
The reduction, formalization, and manuscript were developed with AI assistance.
The reliability theorem and union-estimation method are external contributions.
Independent expert review and novelty assessment remain pending.

\begin{thebibliography}{9}
\bibitem{question}
Antoine Amarilli.
\emph{FPRAS to estimate the probability to get a cyclic subgraph of a
directed graph}. Theoretical Computer Science Stack Exchange, September 26,
2023. \href{https://cstheory.stackexchange.com/q/53359}{Question 53359}.
See also the author's
\href{https://a3nm.net/work/research/questions/\#fpras-for-the-probability-of-obtaining-a-cyclic-subgraph-of-a-directed-graph}{open-question list}.

\bibitem{probgraphs}
Antoine Amarilli, Timothy van Bremen, Octave Gaspard, and Kuldeep S. Meel.
\emph{Approximating Queries on Probabilistic Graphs}.
Logical Methods in Computer Science, 21(4:30), 2025.
\href{https://doi.org/10.46298/lmcs-21(4:30)2025}{doi:10.46298/lmcs-21(4:30)2025}.

\bibitem{ffg}
Weiming Feng, Yucheng Fu, and Heng Guo.
\emph{Approximating two-terminal network reliability}.
\href{https://arxiv.org/abs/2608.02523v2}{arXiv:2608.02523v2}, November 3, 2026.
We use Algorithm S and Theorem 9, with the computational model in Section 2.1.

\bibitem{klm}
Richard M. Karp, Michael Luby, and Neal Madras.
\emph{Monte-Carlo approximation algorithms for enumeration problems}.
Journal of Algorithms, 10(3):429--448, 1989.
\href{https://doi.org/10.1016/0196-6774(89)90038-2}{doi:10.1016/0196-6774(89)90038-2}.
\end{thebibliography}
\end{document}
